\emph{Premise.} Hypothesis (P) of the statement is Premise \claimref{synthesis.premise-positivity-of-values}; every statement below about the values of states is under it, relative to the fixed split of the generators.

\emph{Preliminaries.} By Theorem \claimref{synthesis.signed-directions-and-the-integer-clifford-algebra} the algebra is free on the blades $e_{S}$, the images of the elements of the group $G$ of signed directions, and each blade is invertible with $e_{S}^{-1} = \pm e_{S}$. Inversion is an anti-automorphism of $G$ fixing the sign bit, so its linear extension $\dagger$ is an anti-automorphism of the algebra with $x^{\dagger\dagger} = x$. Since $e_{S}e_{T}^{\dagger} = \pm e_{S \triangle T}$ is a scalar only for $S = T$, where it is $1$, $\langle xy^{\dagger}\rangle_{0} = \sum_{S}x_{S}y_{S}$; this form is symmetric and $\langle xx^{\dagger}\rangle_{0} = \sum x_{S}^{2} \ge 0$. A self-adjoint blade satisfies $e_{S} = e_{S}^{\dagger} = e_{S}^{-1}$, so it has square $1$. Two blades commute or anticommute (Theorem \claimref{synthesis.blade-square-signs-f2-quadratic-form}(ii)). The adjoint, observables, configurations and (P) are those of Theorem \claimref{synthesis.chsh-bound-for-two-subsystems-of-a-clifford-algebra-by-a-rational-certificate}.

\emph{(i).} \emph{The frame.} For each $(p,q) \in \mathcal{E}$ the self-adjoint blades are enumerated (the check does this for all fourteen): they are $1$, a set $F$ of pairwise anticommuting blades, and, in the five signatures named, the pseudoscalar $J$ (central, $J^{\dagger} = J$, $J^{2} = 1$) and the blades $Jh$, $h \in F$. Moreover the products $hk$ ($h \ne k$ in $F$, one for each unordered pair) are blades distinct from each other and from $1$ and the elements of $F$, also after multiplication by $J$. The sizes of $F$ are $0, 0$; $2, 2, 0$; $3, 2, 3, 0$; $5, 5, 5$; $5, 5$ in the order $(1,0), (0,1)$; $(2,0), (1,1), (0,2)$; $(3,0), (2,1), (1,2), (0,3)$; $(4,0), (1,3), (0,4)$; $(5,0), (1,4)$.

\emph{Blocks.} Without $J$ put $P = 1$. With $J$, the elements $P_{\pm} = \tfrac12(1 \pm J)$ are central self-adjoint idempotents with $P_{+} + P_{-} = 1$, $P_{+}P_{-} = 0$, and every statement below is made in the algebra $P_{\pm}\Cl(p,q)$ with unit $P_{\pm}$, where $Jh = \pm h$. In a block $P$, a self-adjoint element is $P(s + v)$ with $s$ a scalar and $v = \sum_{h \in F}v_{h}h$, uniquely; and since the $h$ anticommute and have square $1$, $v^{2} = Q(v) := \sum v_{h}^{2}$ and $vv' - v'v = 2\sum_{h < k}(v_{h}v'_{k} - v_{k}v'_{h})hk$.

\emph{Observables of a block.} $(s + v)^{2} = s^{2} + Q(v) + 2sv$. If this is $1$ in the block then $sv = 0$ (the blades of $F$ are independent of $1$ in the block), so either $v = 0$ and $s = \pm 1$, or $s = 0$ and $Q(v) = 1$. If two observables $v, v'$ with $s = s' = 0$ commute, then all $v_{h}v'_{k} - v_{k}v'_{h}$ vanish (the products $hk$ are independent in the block), so $v$ and $v'$ are proportional, and $Q(v) = Q(v') = 1$ gives $v' = \pm v$.

\emph{Conclusion.} Fix a block and suppose the components of $A_{1}, A_{2}$ in it do not commute. Then both have $s = 0$ and their vectors $v_{1}, v_{2}$ are not proportional. The component of $B_{j}$ commutes with both; if it had $s = 0$ its vector would be $\pm v_{1}$ and $\pm v_{2}$, which is impossible, so it is $\pm P$. Hence the components of $B_{1}, B_{2}$ commute. In every block one of the two commutators vanishes, so $[A_{1},A_{2}][B_{1},B_{2}] = 0$; then $C^{2} = 4$ by clause (i) of Theorem \claimref{synthesis.chsh-bound-for-two-subsystems-of-a-clifford-algebra-by-a-rational-certificate} and, under (P), $\omega(C)^{2} \le \omega(C^{2}) = 4$ by its clause (iii). \emph{Frames.} If $\alpha$ is an automorphism of the algebra carrying one orthonormal frame to another, the adjoint of the second frame is $\alpha \circ \dagger \circ \alpha^{-1}$, and $\alpha$ carries observables and configurations for the first adjoint to those for the second and preserves $[A_{1},A_{2}][B_{1},B_{2}] = 0$; so the statement holds for every frame.

\emph{(ii).} For the thirteen signatures with $4 \le n \le 6$ outside $\mathcal{E}$ the check lists four self-adjoint blades with $A_{1}A_{2} = -A_{2}A_{1}$, $B_{1}B_{2} = -B_{2}B_{1}$ and $A_{i}B_{j} = B_{j}A_{i}$. For $n \ge 7$ choose six of the generators; the subalgebra they generate is $\Cl(p',q')$ with $p' + q' = 6$, which is not in $\mathcal{E}$, its blades are blades of $\Cl(p,q)$ with the same inverses, and so the configuration is a configuration of blades of $\Cl(p,q)$. Every signature with $n \ge 4$ outside the five listed is covered.

The rest of the clause is the blade lemma, clause (iv) of Theorem \claimref{synthesis.chsh-bound-for-two-subsystems-of-a-clifford-algebra-by-a-rational-certificate}, applied to this configuration in $\Cl_{k}(p,q)$: it gives $C^{2} = 4 + 4K$, the functionals $\omega_{r}$, that they satisfy (P) exactly when $8r^{2} \le 1$, and their values $8r$.

\emph{(iii).} By Theorem \claimref{transport.classification-real-clifford-algebras} $\Cl(p,q)$ is $\Mat_{k}(D)$ or a sum of two copies, $D$ one of $\R$, $\C$, $\HH$ according to $p - q \bmod 8$, with $k^{2}\dim D$ (times $2$ for a sum) equal to $2^{n}$. Computing $k$ from the table for each signature gives $k \le 2$ exactly on $\mathcal{E}$ for $n \le 6$ (the check does this up to $n = 8$), and $k \ge 4$ for $n \ge 7$ since $k^{2} \ge 2^{n}/8$. The named cases: $\Cl(3,0) = \Mat_{2}(\C)$, $\Cl(1,3) = \Mat_{2}(\HH)$, $\Cl(3,1) = \Cl(2,2) = \Mat_{4}(\R)$. The label $(p,q)$ is that of the stored theorem ($p$ generators of square $1$). Two commuting copies of $\Cl(1,3)$, or $\Cl(1,3)$ with a further unit commuting with it (which is $\Cl(4,1) = \Mat_{4}(\C)$ by the same table), are algebras outside $\mathcal{E}$, to which (ii) applies.
